\section{Замкнутость КС языков относительно операций}
\tikzsetfigurename{CFG_Closure_}

\begin{theorem}
    Контекстно-свободные языки замкнуты относительно следующих операций:
    \begin{enumerate}
        \item Объединение: если $L_1$ и $L_2$~--- контекстно-свободные языки, то и $L_3 = L_1 \cup L_2$~--- контекстно-свободный.
        \item Конкатенация: если $L_1$ и $L_2$~--- контекстно-свободные языки, то и $L_3 = L_1 \cdot L_2$~--- контекстно-свободный.
        \item Замыкание Клини: если $L_1$~--- контекстно-свободный, то и $L_2 = \bigcup\limits_{i=0}^{\infty} L_1^i $~--- контекстно-свободный.
        \item Разворот: если $L_1$~--- контекстно-свободный, то и $L_2 = {L_1}^r = \{ l^r \mid l \in L_1\}$ является контекстно-свободным.
        \item Пересечение с регулярными языками: если $L_1$~--- контекстно-свободный, а $L_2$~--- регулярный, то  $L_3 = L_1 \cap L_2$~--- контекстно-свободный.
        \item Разность с регулярными языками: если $L_1$~--- контекстно-свободный, а $L_2$~--- регулярный, то  $L_3 = L_1 \setminus L_2$~--- контекстно-свободный.
    \end{enumerate}
\end{theorem}
\begin{proof}
    Для доказательства пунктов 1--4 можно построить КС грамматику нового языка имея грамматики для исходных.
    Будем предполагать, что множества нетерминальных символов различных грамматик для исходных языков не пересекаются.
    \begin{enumerate}
        \item $G_1=\langle\Sigma_1,N_1,P_1,S_1\rangle$~--- грамматика для $L_1$, $G_1=\langle\Sigma_2,N_2,P_2,S_2\rangle$~--- грамматика для $L_2$, тогда
              \[G_3=\langle\Sigma_1 \cup \Sigma_2, N_1 \cup N_2 \cup \{S_3\}, P_1 \cup P_2 \cup \{S_3 \to S_1 \mid S_2\} ,S_3\rangle\]~---~грамматика для $L_3$.
        \item $G_1=\langle\Sigma_1,N_1,P_1,S_1\rangle$~--- грамматика для $L_1$, $G_1=\langle\Sigma_2,N_2,P_2,S_2\rangle$~--- грамматика для $L_2$, тогда $G_3=\langle\Sigma_1 \cup \Sigma_2, N_1 \cup N_2 \cup \{S_3\}, P_1 \cup P_2 \cup \{S_3 \to S_1 S_2\} ,S_3\rangle$~--- грамматика для $L_3$.
        \item $G_1=\langle\Sigma_1,N_1,P_1,S_1\rangle$~--- грамматика для $L_1$, тогда $G_2=\langle\Sigma_1, N_1 \cup \{S_2\}, P_1 \cup \{S_2 \to S_1 S_2\ \mid \varepsilon\}, S_2\rangle$~--- грамматика для $L_2$.
        \item $G_1=\langle\Sigma_1,N_1,P_1,S_1\rangle$~--- грамматика для $L_1$, тогда $G_2=\langle\Sigma_1, N_1, \{N^i \to \omega^R \mid N^i \to \omega \in P_1 \}, S_1\rangle$~--- грамматика для $L_2$.
    \end{enumerate}

    \mytodo{Более вменяеое доказательство замкнутости. Или хотя бы более вменяемая идея. Главное, чтобы не пришлось вводить стековый автомати и прочий ужас.}
    Чтобы доказать замкнутость относительно пересечения с регулярными языками, построим по КС грамматике рекурсивный автомат $R_1$, по регулярному выражению~--- детерминированный конечный автомат $R_2$, и построим их прямое произведение $R_3$.
    Переходы по терминальным символам в новом автомате возможны тогда и только тогда, когда они возможны одновременно и в исходном рекурсивном автомате и в исходном конечном.
    За рекурсивные вызовы отвечает исходный рекурсивный автомат.
    \marginnote{TODO: Вот тут RSM очень внезапно вылез}
    Значит цепочка принимается $R_3$ тогда и только тогда, когда она принимается одновременно $R_1$ и $R_2$: так как состояния $R_3$~--- это пары из состояния $R_1$ и $R_2$, то по трассе вычислений $R_3$ мы всегда можем построить трассу для $R_1$ и $R_2$ и наоборот.

    Чтобы доказать замкнутость относительно разности с регулярным языком, достаточно вспомнить, что регулярные языки замкнуты относительно дополнения, и выразить разность через пересечение с дополнением:
    \[L_1 \setminus L_2 = L_1 \cap \overline{L_2} \qedhere\]
\end{proof}

\begin{theorem}
    Контекстно-свободные языки не замкнуты относительно следующих операций:
    \begin{enumerate}
        \item Пересечение: если $L_1$ и $L_2$~--- контекстно-свободные языки, то и $L_3 = L_1 \cap L_2$~--- не контекстно-свободный.
        \item Разность: если $L_1$ и $L_2$~--- контекстно-свободные языки, то и $L_3 = L_1 \setminus L_2$~--- не контекстно-свободный.
    \end{enumerate}
\end{theorem}

\begin{proof}
    Чтобы доказать незамкнутость относительно пресечения, рассмотрим языки $L_1 = \{a^n b^n c^k \mid n \geq 0, k \geq 0\}$ и $L_2 = \{a^k b^n c^n \mid n \geq 0, k \geq 0\}$.
    Очевидно, что $L_1$ и $L_2$~--- контекстно-свободные языки.
    Рассмотрим $L_3 = L_1 \cap L_2 = \{a^n b^n c^n \mid n \geq 0\}$.
    $L_3$ не является контекстно-свободным по лемме о накачке для контекстно-свободных языков.

    Чтобы доказать незамкнутость относительно разности проделаем следующее.
    \begin{enumerate}
        \item Рассмотрим языки $L_4 = \{a^m b^n c^k \mid m \neq n, k \geq 0\}$ и $L_5 = \{a^m b^n c^k \mid n \neq k, m \geq 0\}$.
              Эти языки являются контекстно-свободными.
              Это легко заметить, если знать, что язык $L'_4 = \{a^m b^n c^k \mid 0 \leq m < n, k \geq 0\}$ задаётся следующей грамматикой:
              \begin{align*}
                  S & \to S c & T & \to a T b \\
                  S & \to T   & T & \to T b   \\
                    &         & T & \to b.
              \end{align*}
        \item Рассмотрим язык $L_6 = \overline{L'_6} = \overline{\{a^n b^m c^k \mid n \geq 0, m \geq 0, k \geq 0\}}$.
              Данный язык является регулярным\sidenote{Предлагаем читателю самостоятельно написать регулярное выражение, задающее этот язык.}.
        \item Рассмотрим язык $L_7 = L_4 \cup L_5 \cup L_6$~--- контекстно-свободный, так как является объединением контекстно-свободных.
        \item Рассмотрим $\overline{L_7} = \{a^n b^n c^n \mid n \geq 0\} = L_3$: $L_4$ и $L_5$ задают языки с правильным порядком символов, но неравным их количеством, $L_6$ задаёт язык с неправильным порядком символов.
              Из предыдущего пункта мы знаем, что $L_3$  не является контекстно-свободным.
              \qedhere
    \end{enumerate}
\end{proof}




%\section{Вопросы и задачи}
%\begin{enumerate}
%  \item Постройте дерево вывода цепочки $w=aababb$ в грамматике $G=\langle\{a,b\},\{S\},\{S\rightarrow \varepsilon \ | \ a \ S \ b \ S \}, S \rangle$.
%  \item Постройте все левосторонние выводы цепочки $w=ababab$ в грамматике $G=\langle\{a,b\},\{S\},\{S\rightarrow \varepsilon \ | \ a \ S \ b \ | S \ S\}, S \rangle$.
%  \item Постройте все правосторонние выводы цепочки $w=ababab$ в грамматике $G=\langle\{a,b\},\{S\},\{S\rightarrow \varepsilon \ | \ a \ S \ b \ | S \ S\}, S \rangle$.
%  \item \label{t1}Постройте все деревья вывода цепочки $w=ababab$ в грамматике $G=\langle\{a,b\},\{S\},\{S\rightarrow \varepsilon \ | \ a \ S \ b \ | S \ S\}, S \rangle$, соответствующие левосторонним выводам.
%  \item \label{t2}Постройте все деревья вывода цепочки $w=ababab$ в грамматике $G=\langle\{a,b\},\{S\},\{S\rightarrow \varepsilon \ | \ a \ S \ b \ | S \ S\}, S \rangle$, соответствующие правосторонним выводам.
%\end{enumerate}
